\documentclass{beamer}

\input{settings.tex}


\title{Sums of Squares Lyapunov Analysis}
\subtitle{Advanced Control}
\author{by Sergei Savin}
\centering
\date{\mydate}



\begin{document}
\maketitle


\begin{frame}{Content}
	\begin{itemize}
		\item Stability as an SoS Problem
		\item Asymptotic Stability as an SoS Problem
		\item Exponential Stability as an SoS Problem
		\item Quadratic Lyapunov Candidate
		\item Non-Polynomial Systems
		\item Stability Check for a Pendulum
		\item Local Stability
		\item S-Procedure
		\item Local Asymptotic Stability via S-Procedure
		\item Literature
	\end{itemize}
\end{frame}


\begin{frame}{Stability as an SoS Problem}
	\begin{flushleft}
		A system $\dot{x}=f(x)$ is stable if there exists $V(x)$ such that:
		\[
		V(x)>0,\qquad V(0)=0,\qquad \dot V(x)\leq0\footnotemark
		\]
		
		\bigskip
		
		\textbf{Notice:}
		\begin{itemize}
			\item The strict inequality is replaced with the bound:
			\[
			V(x)\geq\alpha x^\top x,\qquad \alpha>0.
			\]
			
			\item Since $\dot V(x)
			=
			\frac{\partial V}{\partial x}f(x),$	we can write $-\frac{\partial V}{\partial x}f(x)\geq0$.
		\end{itemize}
		
		\bigskip
		
		We can formulate the stability test as an SoS program:
		\[
		\begin{array}{ll}
			\text{find} & V\\[0.2em]
			\text{s.t.} &
			V(x)-\alpha x^\top x\in\mathrm{SoS},\\
			&
			-\dfrac{\partial V}{\partial x}f(x)\in\mathrm{SoS}.
		\end{array}
		\]
	\end{flushleft}
	
	\footnotetext{\vspace{1.0em}
		The condition $V(x)\geq\alpha x^\top x$, $\alpha>0$ implies $V(x)\to\infty$ as $\|x\|\to\infty$.}
\end{frame}


\begin{frame}{Asymptotic Stability as an SoS Problem}
	\begin{flushleft}
		A system $\dot{x}=f(x)$	is asymptotically stable if there exists $V(x)$ such that:
		\[
		V(x)>0,\qquad V(0)=0,\qquad
		\textcolor{myblue}{\dot V(x)<0}.
		\]
		
		We replace the strict inequality with
		\[
		\dot V(x)\leq-\epsilon x^\top x,
		\qquad \epsilon>0.
		\]
		
		The resulting SoS program is:
		\[
		\begin{array}{ll}
			\text{find} & V\\[0.2em]
			\text{s.t.} &
			V(x)-\alpha x^\top x\in\mathrm{SoS},\\
			&
			-\dfrac{\partial V}{\partial x}f(x)
			\textcolor{myblue}{-\epsilon x^\top x}\in\mathrm{SoS}.
		\end{array}
		\]
	\end{flushleft}
	
\end{frame}


\begin{frame}{Exponential Stability as an SoS Problem}
	\begin{flushleft}
		A system $\dot{x}=f(x)$	is exponentially stable if there exists $V(x)$ and
		$\alpha,\beta,\epsilon>0$ such that:
		\[
		\alpha x^\top x
		\leq V(x)
		\textcolor{myblue}{\leq \beta x^\top x},\qquad 
		\dot V(x)\leq-\epsilon x^\top x.
		\]
		
		\bigskip
		
		The corresponding SoS program is:
		\[
		\begin{array}{ll}
			\text{find} & V\\[0.2em]
			\text{s.t.} &
			V(x)-\alpha x^\top x\in\mathrm{SoS},\\
			&
			\textcolor{myblue}{-V(x)+\beta x^\top x\in\mathrm{SoS}},\\
			&
			-\dfrac{\partial V}{\partial x}f(x)
			-\epsilon x^\top x\in\mathrm{SoS}.
		\end{array}
		\]
	\end{flushleft}
	
\end{frame}


\begin{frame}{Quadratic Lyapunov Candidate}
	\begin{flushleft}
		We can simplify the process by limiting ourselves to
		quadratic Lyapunov functions:
		\[
		V(x)=\frac{1}{2}x^\top Px,
		\qquad P=P^\top.
		\]
		
		For example, the stability test becomes:
		\[
		\begin{array}{ll}
			\text{find} & P\\[0.2em]
			\text{s.t.} &
			P\succeq\alpha I,\\
			&
			-x^\top Pf(x)\in\mathrm{SoS}.
		\end{array}
		\]
		
		
		Replacing the last condition with:
		\[
		-x^\top Pf(x)-\epsilon x^\top x\in\mathrm{SoS}
		\]
		verifies asymptotic and exponential stability. Indeed, for any $P\succ0$:
		\[
		\lambda_{\max}(P)x^\top x
		\geq V(x)
		\geq
		\lambda_{\min}(P)x^\top x.
		\]
	\end{flushleft}
\end{frame}

\begin{frame}{Non-Polynomial Systems}
	\begin{flushleft}
		Many dynamical systems are non-polynomial.
		For example, mechanical systems often include
		trigonometric functions.
		
		\bigskip
		
		Consider a pendulum:
		\[
		ml^2\ddot{\theta}
		+b\dot{\theta}
		+mgl\sin(\theta)=0.
		\]
		
		\bigskip
		
		Working with such systems requires a polynomial reformulation:
		
		\begin{itemize}
			\item Approximate the nonlinearities using, for example,
			\[
			\sin(x)\approx
			x-\frac{x^3}{3!}+\frac{x^5}{5!}.
			\]
			
			\item Find an equivalent polynomial formulation,
			possibly using additional variables and constraints:
			\[
			s=\sin(x),\qquad
			c=\cos(x),\qquad
			s^2+c^2=1.
			\]
		\end{itemize}
	\end{flushleft}
	
\end{frame}


\begin{frame}{Stability Check for a Pendulum, Part 1}
	\begin{flushleft}
		Consider a simple pendulum with all parameters equal to $1$:
		\[
		\ddot{\theta}+\dot{\theta}+\sin(\theta)=0.
		\]
		
		\bigskip
		
		Define new variables:
		\[
		x_1=\dot{\theta},
		\qquad
		x_2=\sin(\theta),
		\qquad
		x_3=\cos(\theta)-1.
		\]
		
		\bigskip
		
		This change of variables accomplishes four things:
		
		\begin{enumerate}
			\item We transform the second-order ODE into a first-order system.
			
			\item We eliminate the trigonometric functions from the dynamics.
			
			\item The equilibrium of interest becomes $x=0$.
			
			\item We eliminate $\theta$ as a state variable, thus we do not need to establish connection between  $\theta$ and $x_2$ and $x_3$.
		\end{enumerate}
	\end{flushleft}
	
\end{frame}


\begin{frame}{Stability Check for a Pendulum, Part 2}
	\begin{flushleft}
		In the new variables, the dynamics $\ddot{\theta}+\dot{\theta}+\sin(\theta)=0$ become
\[
\begin{aligned}
	\dot{x}_1 &= -x_1-x_2, \\
	\dot{x}_2 &= \cos(\theta)\dot{\theta} = x_1(x_3+1), \qquad  \textcolor{mygray}{\cos(\theta)=x_3+1} \\
	\dot{x}_3 &= -\sin(\theta)\dot{\theta} = -x_1x_2,  \qquad \quad  \textcolor{mygray}{\sin(\theta)=x_2}
\end{aligned}
\]

		
		\bigskip
		
		The relation between sine and cosine is encoded by
		\[
		\sin^2(\theta)+\cos^2(\theta)=1.
		\]
		
		Using the new variables, this becomes
		\[
		h(x)
		=
		x_2^2+(x_3+1)^2-1
		=0.
		\]
		
		\bigskip
		
		Thus, we obtain a polynomial system with a polynomial constraint:
		\[
		\boxed{
			\dot{x}=f(x),
			\qquad
			h(x)=0.
		}
		\]
	\end{flushleft}
\end{frame}


\begin{frame}{Stability Check for a Pendulum, Part 3}
	\begin{flushleft}
		We can formulate the stability check using SoS programming. Importantly, the stability conditions only need to hold on the manifold
		\[
		h(x)=0.
		\]
		
		
		We can use polynomial multipliers to enforce this:
		\[
		\begin{array}{ll}
			\text{find} & V(x),\ \lambda(x),\ \mu(x)\\[0.2em]
			\text{s.t.}
			&
			V(x)-\alpha x^\top x+\mu(x)h(x)
			\in\mathrm{SoS},
			\\[0.4em]
			&
			-\dfrac{\partial V}{\partial x}f(x)
			-\epsilon x^\top x+\lambda(x)h(x)
			\in\mathrm{SoS}.
		\end{array}
		\]
		
		\bigskip
		
		Note that $\lambda(x)$ and $\mu(x)$ are arbitrary polynomial multipliers;
		they are not required to be nonnegative. On the manifold $h(x)=0$, these conditions reduce to
		\[
		V(x)\geq\alpha x^\top x,
		\qquad
		\dot V(x)\leq-\epsilon x^\top x.
		\]
		

	\end{flushleft}
\end{frame}


\begin{frame}{Local Stability}
	\begin{flushleft}
		Many systems, including the pendulum, are only locally stable. To verify local stability, we require the Lyapunov conditions to hold only on a specific region $\mathcal{D}$.
		
		\bigskip
		
		Let us describe $\mathcal{D}$ as a level set of a
		positive-definite function $W(x)$:
		\[
		\mathcal{D}
		=
		\left\{
		x: W(x)\leq c
		\right\},
		\qquad c>0.
		\]
		
		\bigskip
		
		For example, if
		\[
		W(x)=x^\top Px,
		\qquad P\succ0,
		\]
		then
		\[
		\mathcal{D}
		=
		\left\{
		x:x^\top Px\leq c
		\right\}
		\]
		is an ellipsoid.
	\end{flushleft}
\end{frame}


\begin{frame}{S-Procedure}
	\begin{flushleft}
		Suppose we want to prove that a polynomial $f(x)$ is
		nonnegative on the set
		\[
		g(x)\leq0.
		\]
		
		We can search for a polynomial multiplier
		$\lambda(x)\in\mathrm{SoS}$ such that:
		
		\begin{align}
			\label{eq:s-procedure}
			f(x)+\lambda(x)g(x)\in\mathrm{SoS}.
		\end{align}
		
		
		\begin{itemize}
			\item Equation \eqref{eq:s-procedure} implies $
			f(x)+\lambda(x)g(x)\geq0$. For such $x$ that $g(x)\leq0$, $f(x)
			\geq
			-\lambda(x)g(x)
			\geq0.$
			
			\item For such $x$ that $g(x)>0$, however, $-\lambda(x)g(x)\leq0$, 
			so the condition does not imply that $f(x)\geq0$.
		\end{itemize}
		
		
		\bigskip
		
		This allows the SoS program to enforce $f(x)\geq0$
		only on the desired set $g(x)\leq0$.
	\end{flushleft}
\end{frame}




\begin{frame}{Local Asymptotic Stability via S-Procedure}
	\begin{flushleft}
		Using the S-procedure, we can restrict the positivity
		requirement on $V(x)$ and the negativity requirement on
		$\dot V(x)$ to a region $\mathcal{D}$. This allows us to prove \textbf{local asymptotic stability}.
		
		\bigskip
		
		A natural choice of $\mathcal{D}$ is a sublevel set of $V(x)$:
		\[
		\mathcal{D}
		=
		\left\{
		x:V(x)\leq c
		\right\}.
		\]
		
		\bigskip
		
		The SoS formulation becomes
		\[
		\begin{array}{ll}
			\text{find} &
			V,\ \lambda_1,\ \lambda_2
			\\[0.2em]
			\text{s.t.}
			&
			V(x)-\alpha x^\top x
			+\lambda_1(x)\bigl(V(x)-c\bigr)
			\in\mathrm{SoS},
			\\[0.4em]
			&
			-\dfrac{\partial V}{\partial x}f(x)
			-\epsilon x^\top x
			+\lambda_2(x)\bigl(V(x)-c\bigr)
			\in\mathrm{SoS},
			\\[0.4em]
			&
			\lambda_1(x)\in\mathrm{SoS},
			\qquad
			\lambda_2(x)\in\mathrm{SoS}.
		\end{array}
		\]
		{\small
		$\lambda_1$ and $\lambda_2$ are required to be SoS
		(nonnegative) multipliers. The parameters $\alpha>0$
		and $\epsilon>0$ enforce strict Lyapunov inequalities.}
	\end{flushleft}
\end{frame}




\begin{frame}{Literature}

\begin{itemize}
	
	\item Stephen Lall -- Stanford University:
	\bref{https://lall.stanford.edu/ee464/lectures/sum_of_squares.pdf}
	{Sum of Squares.}
	
		\item Russ Tedrake -- MIT:
	\bref{https://underactuated.csail.mit.edu/lyapunov.html}
	{Underactuated Robotics: Lyapunov Analysis, ch. 9.2.2, 9.2.6}
	
	\item MIT OpenCourseWare -- Algebraic Techniques and Semidefinite Optimization:
	\bref{https://ocw.mit.edu/courses/6-972-algebraic-techniques-and-semidefinite-optimization-spring-2006/resources/lecture_10/}
	{Lecture 10: Sums of Squares and Semidefinite Programming.}
	
	\item Antonis Papachristodoulou and Stephen Prajna:
	\bref{https://users.ox.ac.uk/~engs0587/PapPACC05.pdf}
	{A Tutorial on Sum of Squares Techniques for Systems Analysis.}
	
	
\end{itemize}

\end{frame}



\myqrframe




\end{document}
